\documentclass[11pt,oneside]{article}
\usepackage[utf8]{inputenc}
\usepackage[T1]{fontenc}
\usepackage{amsmath, amsthm, amssymb, amsfonts}
\usepackage{geometry}
\usepackage{hyperref}

\geometry{margin=1in}

\title{\textbf{The Soddy-Einstein Gasket: \\ A Strictly Finitist Geometric Foundation for Lazy Constructive Analysis}}
\author{Two Wings}
\date{April 2026}

\begin{document}

\maketitle

\begin{abstract}
The Soddy-Einstein Gasket (SEG) is a geometric framework in which addition corresponds to tangency of spheres with integer curvatures and multiplication is realized as a lazy, pruned pathfinding process. Proofs are closed cyclic tangent paths. 

The gasket is an open recursive process; no part of it is postulated as completed. Numerical objects exist only as finite-depth spheres or finite tangent paths that can be lazily extended on demand. There are no completed real numbers and no completed infinite sets of any kind. 

An optional Axiom of Information Equivalence (\( |P(\mathbb{N})| = |\mathbb{N}| \)) may be adjoined solely to block classical Cantorian diagonalization. Only when this axiom is adopted does the limit set of the gasket become fully addressable, effectively treating the entire fractal as informationally complete in a self-similar sense. Without it, the system remains purely finitist with no completed infinities whatsoever. This yields a decidable base layer while supporting Bishop-style constructive analysis through lazy finite-precision approximations.
\end{abstract}

\section{Introduction}
All mathematical activity in this framework occurs inside an open, recursively generated Soddy (Apollonian) sphere packing. The gasket is never treated as a completed totality — it is a process of successive tangency governed by Descartes' Circle Theorem. At every stage one works only with finite-depth spheres and finite tangent paths. 

No completed infinite sets are assumed. The natural numbers, the power set, and the continuum are never reified as finished objects.

\section{The Metric-Cardinality Postulate}
Information capacity \(\mathcal{C}(s)\) of any object \(s\) (a sphere in the current finite-depth packing) is defined as the inverse of its curvature:

\begin{equation}
\mathcal{C}(s) = \frac{1}{\kappa(s)}
\end{equation}

Curvature \(\kappa = 0\) at the origin corresponds to unbounded potential. Any specific definition or constant introduces \(\kappa > 0\), locally bounding information while preserving the recursive structure.

\section{Base Geometry and Operations}
Addition is realized directly as tangency between spheres. Integer curvatures are preserved under Descartes' Circle Theorem, keeping the structure discrete and decidable at each finite depth.

Multiplication \(x \cdot y\) is implemented as a \textbf{lazy pruned navigation}:
\begin{itemize}
    \item Start from the spheres representing \(x\) and \(y\).
    \item Generate candidate product spheres by descending deeper into the gasket.
    \item Prune the search using Presburger-style linear constraints together with divisibility, parity, and magnitude bounds.
    \item Compute only to the exact finite depth required by the query or desired precision.
\end{itemize}

This keeps the base layer Presburger-like (decidable) while allowing non-linear operations through explicit, terminating computation.

\section{Proof Theory: Cyclic Tangent Paths}
A formal proof is a finite closed cycle of tangent steps that returns to the starting sphere.

\begin{itemize}
    \item Consistency is geometric: incompatible curvatures cannot occupy the same location.
    \item Every inference corresponds to a verified tangency.
\end{itemize}

No completed objects are required for verification — only the finite path.

\section{Optional Axiom of Information Equivalence}
The system functions completely without any axiom concerning power sets or cardinalities. 

For the sole purpose of countering non-constructive Cantorian diagonalization, one may optionally introduce:
\begin{equation}
|P(\mathbb{N})| = |\mathbb{N}|
\end{equation}

Only when this axiom is adopted does the limit set (“Cantor dust”) of the gasket become fully addressable via self-similarity. In that case — and only then — one can regard the entire fractal as informationally complete in a self-similar sense. Without the axiom, the gasket remains an open process with no completed infinite aspect whatsoever.

\section{Constructive Approximations (Strictly No Completed Reals)}
There are no real numbers as completed objects. Any numerical quantity is represented by a finite-depth sphere whose curvature bounds the uncertainty. 

Higher precision is obtained by lazy descent to a deeper tangent child sphere using the pruned multiplication engine. The process always terminates for any explicitly requested finite precision, thanks to effective geometric bounds and pruning rules.

This framework supports Bishop-style constructive analysis (derivatives, integrals, solutions to equations, etc.) as explicit transformations of tangent paths, executed lazily to any demanded modulus of continuity. The “continuum” is never actualized — it exists only as potential deeper navigation within the open gasket.

\section{Inversion of Diagonalization (Under the Optional Axiom)}
When \( |P(\mathbb{N})| = |\mathbb{N}| \) is adopted, Cantor’s diagonal argument is inverted. Any purported uncountable real requires an index \( I \) whose digit cannot be reached by any path in the gasket. Self-similarity plus the axiom ensures no such unreachable index exists; the diagonal number is merely a deeper branch already indexed within the structure.

Without the axiom the classical argument does not apply, because no completed power set or continuum is ever postulated.

\section{Conclusion}
The Soddy-Einstein Gasket is a strictly finitist geometric foundation. The gasket itself is an open recursive process, never postulated as completed. All objects and proofs remain finite or lazily extensible. Addition is tangency, multiplication is demand-driven pruned navigation, and proofs are cyclic tangent paths.

Only the optional Axiom of Information Equivalence introduces a completed infinite aspect (by making the limit set fully addressable). Without it, the system contains no completed infinities at all. This yields a clean, decidable, and geometrically consistent basis for lazy constructive analysis while avoiding non-constructive commitments.

\end{document}


OLDER

\documentclass[11pt,oneside]{article}
\usepackage[utf8]{inputenc}
\usepackage[T1]{fontenc}
\usepackage{amsmath, amsthm, amssymb, amsfonts}
\usepackage{geometry}
\usepackage{hyperref}
\usepackage{cite}

\geometry{margin=1in}

\title{\textbf{The Soddy-Einstein Gasket: \\ A Finitist Metric-Cardinality Framework for Constructive Calculus}}
\author{Two Wings}
\date{April 2026}

\begin{document}

\maketitle

\begin{abstract}
We introduce the Soddy-Einstein Gasket (SEG), a geometric foundation for mathematics that replaces the completed infinite set \(\mathbb{N}\) with recursive sphere packings governed by Descartes' Circle Theorem. Cardinality is defined as the inverse of curvature in an Apollonian/Soddy sphere packing. Addition corresponds to tangency, multiplication to a lazy pruned navigation of the self-similar structure, and proofs to closed cyclic tangent paths. 

A single Axiom of Information Equivalence, \( |P(\mathbb{N})| = |\mathbb{N}| \), collapses the alleged gap between countable and uncountable, making the limit set (Cantor dust) addressable via infinite-depth but finitely describable paths. This yields a decidable base layer (Presburger-like) while supporting a full Bishop-style constructive calculus through lazy computation. We invert Cantor's diagonal argument by demanding a witness digit unreachable by the gasket's enumeration.
\end{abstract}

\section{Introduction}
Standard foundations separate discrete arithmetic from continuous analysis and struggle with the tension between decidability (Presburger arithmetic) and expressive power (Peano arithmetic with multiplication). We propose a unified geometric manifold—the Soddy Gasket—where logical objects are spheres with integer curvatures, and information flow follows physical-like constraints inspired by the Einstein Equivalence Principle.

In this framework, ``infinity'' is never completed; it is the unbounded depth of recursive subdivision. The system remains finitist in its processes while allowing arbitrarily precise constructive mathematics.

\section{The Metric-Cardinality Postulate}
We postulate that the information capacity \(\mathcal{C}(s)\) of a logical object \(s\) (a sphere in the gasket) is inversely proportional to its curvature \(\kappa(s)\):

\begin{equation}
\mathcal{C}(s) = \frac{1}{\kappa(s)}
\label{eq:metric}
\end{equation}

- At the origin (identity, flat vacuum, \( |x| = 0 \)): \(\kappa = 0\), so \(\mathcal{C} \to \infty\). This is the potential for all paths.
- For any concrete object (\( |x| > 0 \)): \(\kappa > 0\), bounding local information and enforcing decidability at each finite depth.

This mirrors the Equivalence Principle: curvature (constraint) locally limits information density, while the global structure preserves infinite capacity through self-similarity.

\section{Axioms and Base Theory}
The foundation is the first-order theory of addition interpreted geometrically as tangency operations in the Soddy gasket (governed by Descartes' Circle Theorem for integer curvatures).

We introduce exactly one non-Presburger axiom:

\begin{equation}
|P(\mathbb{N})| = |\mathbb{N}|
\label{eq:info}
\end{equation}

This asserts that the power set of paths through the gasket has the same cardinality as the set of finite paths. The ``Cantor dust'' (limit set) is not an extra uncountable layer but is bijectively indexed by (potentially infinite-depth) paths in the self-similar structure. There is no completed actual infinity; only recursive growth.

\section{Proof Theory: Cyclic Tangent Paths}
A proof is not a linear deduction but a \textbf{closed geodesic} (cyclic tangent path) in the tangency graph of the gasket.

\begin{itemize}
    \item \textbf{Soundness and Consistency}: A contradiction \( P \land \neg P \) would require two spheres of incompatible curvatures to occupy the same location, violating the packing condition (non-overlapping interiors).
    \item \textbf{Verification}: A statement holds if there exists a finite sequence of tangency steps (lazy additions) that forms a closed loop returning to the starting sphere. This corresponds to a self-consistent identity in the geometry.
\end{itemize}

This topological notion of proof aligns with constructive mathematics: existence is witnessed by a concrete path.

\section{Multiplication as Lazy Navigation}
Multiplication \( x \cdot y = z \) is not a primitive relation but a \textbf{lazy recursive pathfinding} subroutine:

\begin{enumerate}
    \item Start from spheres representing \( x \) and \( y \).
    \item Use geometric constraints (Descartes' theorem and bounds from Presburger addition) to prune the search space of candidate spheres for \( z \).
    \item Enumerate candidate tangency paths only to the finite depth required by the current precision or query.
    \item The ``For Any'' quantifier is resolved by explicit construction of a witness path; no universal quantification over a completed set is needed.
\end{enumerate}

This keeps the base Presburger layer decidable (linear constraints) while allowing non-linear operations through bounded computation. The pruning uses algebraic properties such as divisibility, parity, and magnitude bounds—exactly as in modern SMT solvers.

\section{Constructive Calculus in the Gasket}
Real numbers are represented not as completed Cauchy sequences but as \textbf{lazy convergent sequences of curvatures} \(\{ \kappa_n \}\) where each subsequent sphere is tangent to the previous, with error controlled by \( 1/\kappa_n \).

Multiplication and other operations on reals become composition of path-generators, executed lazily to any requested modulus of continuity (Bishop-style).

The Intermediate Value Theorem, for example, is realized via a constructive bisection-like search along tangent paths, guaranteed to terminate at any finite precision because the gasket provides effective bounds.

\section{Inversion of Cantor's Diagonal Argument}
We turn Cantor's diagonalization against the claim of uncountability. Suppose a real number \( r = 0.d_1 d_2 d_3 \dots \) in \([0,1]\) is asserted to lie outside any enumeration of countable paths.

The proponent must exhibit a specific index \( I \) such that digit \( d_I \) cannot be computed by any path in the Soddy-Einstein Gasket. However:

- The gasket is self-similar and space-filling under Axiom \eqref{eq:info}.
- Every possible finite prefix of digits corresponds to a finite-depth sphere.
- Every infinite sequence is an infinite-depth path whose address is indexed within the structure (since \( |P(\mathbb{N})| = |\mathbb{N}| \)).

Thus, no such unreachable index \( I \) can be provided. The ``uncountable'' reduces to the ``not yet computed to sufficient depth.'' Existence in the continuum is equivalent to addressability via a (lazy) path in the gasket.

\section{Conclusion}
The Soddy-Einstein Gasket offers a finitist, geometric reconstruction of mathematics. Presburger arithmetic provides the coarse, decidable geometry of large spheres. Lazy multiplication supplies the fine navigation needed for analysis. Cyclic paths give a natural proof theory, and the single Axiom of Information Equivalence eliminates the ontological excess of transfinite set theory while preserving the power of constructive calculus.

In this framework, incompleteness manifests only as measure-zero limit points that do not obstruct practical computation or proof at any finite precision. Mathematics becomes navigation of a single, self-similar, computable crystal.

\end{document}